\documentclass[10pt,a4paper]{amsart}
\usepackage{amsmath,amssymb,amsthm,graphicx}
% Retained preamble for the Stacks Project chapter snapshot in this directory.
% The chapter's own preamble.tex is a copy of the shared Stacks Project preamble; the block
% below is sliced MECHANICALLY, line for line, from it. Nothing was added, removed or
% redefined; the wrapper class, the xr-hyper/hyperref wiring and the externaldocument list
% are omitted because the assembler supplies the document class.
%
% lines taken from preamble.tex: 3, 5, 9, 11, 12, 156, 157, 158, ... (51 lines total)
%
% package coverage restored from the chapter's own preamble (2026-10-13): xy, xy-2cell

\usepackage[all]{xy}
\xyoption{2cell}
\UseAllTwocells

\newenvironment{reference}{\comment}{\endcomment}
\newenvironment{slogan}{\comment}{\endcomment}
\newenvironment{history}{\comment}{\endcomment}
\theoremstyle{plain}
\newtheorem{theorem}[subsection]{Theorem}
\newtheorem{proposition}[subsection]{Proposition}
\newtheorem{lemma}[subsection]{Lemma}
\theoremstyle{definition}
\newtheorem{definition}[subsection]{Definition}
\newtheorem{example}[subsection]{Example}
\newtheorem{exercise}[subsection]{Exercise}
\newtheorem{situation}[subsection]{Situation}
\theoremstyle{remark}
\newtheorem{remark}[subsection]{Remark}
\newtheorem{remarks}[subsection]{Remarks}
\numberwithin{equation}{subsection}
\def\lim{\mathop{\mathrm{lim}}\nolimits}
\def\colim{\mathop{\mathrm{colim}}\nolimits}
\def\Spec{\mathop{\mathrm{Spec}}}
\def\Hom{\mathop{\mathrm{Hom}}\nolimits}
\def\Ext{\mathop{\mathrm{Ext}}\nolimits}
\def\SheafHom{\mathop{\mathcal{H}\!\mathit{om}}\nolimits}
\def\SheafExt{\mathop{\mathcal{E}\!\mathit{xt}}\nolimits}
\def\Sch{\mathit{Sch}}
\def\Mor{\mathop{\mathrm{Mor}}\nolimits}
\def\Ob{\mathop{\mathrm{Ob}}\nolimits}
\def\Sh{\mathop{\mathit{Sh}}\nolimits}
\def\NL{\mathop{N\!L}\nolimits}
\def\CH{\mathop{\mathrm{CH}}\nolimits}
\def\proetale{{pro\text{-}\acute{e}tale}}
\def\etale{{\acute{e}tale}}
\def\QCoh{\mathit{QCoh}}
\def\Ker{\mathop{\mathrm{Ker}}}
\def\Im{\mathop{\mathrm{Im}}}
\def\Coker{\mathop{\mathrm{Coker}}}
\def\Coim{\mathop{\mathrm{Coim}}}
\DeclareMathSymbol{\boxtimes}{\mathbin}{AMSa}{"02}
\def\QCohstack{\mathcal{QC}\!\mathit{oh}}
\def\Cohstack{\mathcal{C}\!\mathit{oh}}
\def\Spacesstack{\mathcal{S}\!\mathit{paces}}
\def\Quotfunctor{\mathrm{Quot}}
\def\Hilbfunctor{\mathrm{Hilb}}
\def\Curvesstack{\mathcal{C}\!\mathit{urves}}
\def\Polarizedstack{\mathcal{P}\!\mathit{olarized}}
\def\Complexesstack{\mathcal{C}\!\mathit{omplexes}}
\def\Pic{\mathop{\mathrm{Pic}}\nolimits}
\def\Picardstack{\mathcal{P}\!\mathit{ic}}
\def\Picardfunctor{\mathrm{Pic}}
\def\Deformationcategory{\mathcal{D}\!\mathit{ef}}
\usepackage{cleveref}
\crefname{originallemma}{original lemma}{original lemmas}
\Crefname{originallemma}{Original lemma}{Original lemmas}
\makeatletter
\newlabel{more-algebra-lemma-valuation-ring-torsion-free-flat}{{0539}{source}{}{sourcepointer-1}{}}
\newlabel{more-algebra-lemma-valuation-ring-torsion-free-flat@cref}{{[originallemma][0][]0539}{[1][0][]source}{}{}{}}
\makeatother
\crefname{originallemma}{original lemma}{original lemmas}
\Crefname{originallemma}{Original lemma}{Original lemmas}
\makeatletter
\newlabel{spaces-flat-theorem-finite-type-flat}{{0DLR}{source}{}{sourcepointer-2}{}}
\newlabel{spaces-flat-theorem-finite-type-flat@cref}{{[originallemma][0][]0DLR}{[1][0][]source}{}{}{}}
\makeatother
\crefname{originallemma}{original lemma}{original lemmas}
\Crefname{originallemma}{Original lemma}{Original lemmas}
\makeatletter
\newlabel{spaces-morphisms-definition-valuative-criterion}{{03IX}{source}{}{sourcepointer-3}{}}
\newlabel{spaces-morphisms-definition-valuative-criterion@cref}{{[originallemma][0][]03IX}{[1][0][]source}{}{}{}}
\makeatother
\crefname{originallemma}{original lemma}{original lemmas}
\Crefname{originallemma}{Original lemma}{Original lemmas}
\makeatletter
\newlabel{spaces-properties-proposition-quasi-coherent}{{03M3}{source}{}{sourcepointer-4}{}}
\newlabel{spaces-properties-proposition-quasi-coherent@cref}{{[originallemma][0][]03M3}{[1][0][]source}{}{}{}}
\makeatother
\title{Valuative existence for the Quot functor}
\author{The Stacks Project Authors}
\date{stacks-project:moduli.tex:local-snapshot-20261012; GNU FDL 1.2 or later}
\begin{document}
\maketitle
\noindent\textit{Editorial note.} Statement span is lines 769-776 of sources/stacks-moduli/pinned-original-source.tex; proof span is lines 778-825, the complete author proof (46 body lines).\par\medskip
\noindent\textit{Editorial note.} Navigation stubs: 4 cited-but-not-included labels are supplied as non-load-bearing pointers with their pinned tags where the snapshot tag index has them.\par\medskip
\noindent\textit{Editorial note.} extra\_packages is empty on purpose: the retained preamble already loads xy with options and a second declaration clashes in the assembler wrapper.\par\medskip
\noindent\textit{Editorial note.} Text is retained verbatim, including the author's own typos and placeholders.\par\medskip
\noindent\textit{Editorial note.} Queue index 6 of queue\_moduli.json (pb=2311, stmt 769-776, proof 778-825).\par\medskip
\noindent\textit{Editorial note.} stacks-moduli wave-1: fresh adjudication accept (2026-10-14); pool per Chief order 08:5x (dpa superseded to lane-s9); moduli.tex 53 labels/48 with proof judged 25 accept / 23 reject / 5 pb=0 env-dup rows excluded; spans cite-free and reference/slogan/history-free so no patchenv needed.\par\medskip
\noindent\textit{Allowed standard prerequisites (editorial).}\begin{itemize}
\item Explicit prerequisite boundary. The cited inputs are prerequisites; the argument in the retained span is not replaced by any cited result.
\end{itemize}\medskip
\ifdefined\hypertarget\hypertarget{sourcepointer-1}{}\fi
\noindent\textit{Source reference.} Original lemma 0539: sources/stacks-moduli/pinned-original-source.tex, source lines 825--825. Cited by the retained statement or proof but not part of the collected spans; supplied only as a navigation pointer. The argument does not depend on its content.\par\medskip
\ifdefined\hypertarget\hypertarget{sourcepointer-2}{}\fi
\noindent\textit{Source reference.} Original lemma 0DLR: sources/stacks-moduli/pinned-original-source.tex, source lines 825--825. Cited by the retained statement or proof but not part of the collected spans; supplied only as a navigation pointer. The argument does not depend on its content.\par\medskip
\ifdefined\hypertarget\hypertarget{sourcepointer-3}{}\fi
\noindent\textit{Source reference.} Original lemma 03IX: sources/stacks-moduli/pinned-original-source.tex, source lines 825--825. Cited by the retained statement or proof but not part of the collected spans; supplied only as a navigation pointer. The argument does not depend on its content.\par\medskip
\ifdefined\hypertarget\hypertarget{sourcepointer-4}{}\fi
\noindent\textit{Source reference.} Original lemma 03M3: sources/stacks-moduli/pinned-original-source.tex, source lines 825--825. Cited by the retained statement or proof but not part of the collected spans; supplied only as a navigation pointer. The argument does not depend on its content.\par\medskip
\begin{lemma}
\label{lemma-quot-existence-part}
Assume $X \to B$ is proper as well as of finite presentation
and $\mathcal{F}$ quasi-coherent of finite type.
Then $\Quotfunctor_{\mathcal{F}/X/B} \to B$ satisfies the existence part
of the valuative criterion (Morphisms of Spaces, Definition
\ref{spaces-morphisms-definition-valuative-criterion}).
\end{lemma}

\begin{proof}
Taking base change, this immediately reduces to the following
problem: given a valuation ring $R$ with fraction field $K$,
an algebraic space $X$ proper over $R$, a finite type quasi-coherent
$\mathcal{O}_X$-module $\mathcal{F}$, and a coherent
quotient $\mathcal{F}_K \to \mathcal{Q}_K$, show there exists
a quotient $\mathcal{F} \to \mathcal{Q}$ where $\mathcal{Q}$ is a
finitely presented $\mathcal{O}_X$-module
flat over $R$ whose generic fibre is $\mathcal{Q}_K$.
Observe that by Flatness on Spaces, Theorem
\ref{spaces-flat-theorem-finite-type-flat}
any finite type quasi-coherent $\mathcal{O}_X$-module
$\mathcal{F}$ flat over $R$ is of finite presentation.
We first solve the existence of $\mathcal{Q}$ affine locally.

\medskip\noindent
Affine locally we arrive at the following problem:
let $R \to A$ be a finitely presented ring map,
let $M$ be a finite $A$-module, let $\varphi : M_K \to N_K$ be
an $A_K$-quotient module. Then we may consider
$$
L = \{x \in M \mid \varphi(x \otimes 1) = 0 \}
$$
The $M \to M/L$ is an $A$-module quotient which is
torsion free as an $R$-module. Hence it is flat as an
$R$-module (More on Algebra, Lemma
\ref{more-algebra-lemma-valuation-ring-torsion-free-flat}).
Since $M$ is finite as an $A$-module so is $L$ and we
conclude that $L$ is of finite presentation as an $A$-module
(by the reference above). Clearly $M/L$ is the unique such
quotient with $(M/L)_K = N_K$.

\medskip\noindent
The uniqueness in the construction of the previous paragraph
guarantees these quotients glue and give the desired $\mathcal{Q}$.
Here is a bit more detail. Choose a surjective \'etale morphism
$U \to X$ where $U$ is an affine scheme. Use the above construction
to construct a quotient $\mathcal{F}|_U \to \mathcal{Q}_U$
which is quasi-coherent, is flat over $R$, and recovers $\mathcal{Q}_K|U$
on the generic fibre. Since $X$ is separated, we see that
$U \times_X U$ is an affine scheme \'etale over $X$ as well.
Then $\mathcal{F}|_{U \times_X U} \to \text{pr}_1^*\mathcal{Q}_U$ and
$\mathcal{F}|_{U \times_X U} \to \text{pr}_2^*\mathcal{Q}_U$
agree as quotients by the uniquess in the construction. Hence we may descend
$\mathcal{F}|_U \to \mathcal{Q}_U$ to a surjection
$\mathcal{F} \to \mathcal{Q}$ as desired (Properties of Spaces,
Proposition \ref{spaces-properties-proposition-quasi-coherent}).
\end{proof}

\end{document}
